\documentclass[11pt]{article}

\usepackage[a4paper,margin=1.05in]{geometry}
\usepackage{amsmath,amssymb,amsthm}
\usepackage{mathtools}
\usepackage{stmaryrd}
\usepackage[protrusion=true,expansion=false]{microtype}
\usepackage{enumitem}
\usepackage[dvipsnames]{xcolor}
\usepackage{booktabs}
\usepackage[hidelinks]{hyperref}

\hypersetup{
  colorlinks=true,
  linkcolor=BrickRed,
  citecolor=BrickRed,
  urlcolor=BrickRed
}

\theoremstyle{plain}
\newtheorem{theorem}{Theorem}
\newtheorem{proposition}[theorem]{Proposition}
\newtheorem{corollary}[theorem]{Corollary}
\newtheorem{conjecture}[theorem]{Conjecture}
\newtheorem{observation}[theorem]{Observation}
\theoremstyle{definition}
\newtheorem{definition}[theorem]{Definition}
\newtheorem{example}[theorem]{Example}
\newtheorem{obligation}{Obligation}
\theoremstyle{remark}
\newtheorem{remark}[theorem]{Remark}

\newcommand{\Lam}{\Lambda}
\newcommand{\rhoc}{\textsf{rho}}
\newcommand{\bis}{\;\approx\;}
\newcommand{\bisW}[1]{\approx_{#1}}
\newcommand{\treq}[1]{\simeq^{\mathrm{tr}}_{#1}}
\newcommand{\inp}[3]{\mathsf{for}(#1 \leftarrow #2)\,#3}
\newcommand{\out}[2]{#1!(#2)}
\newcommand{\quo}[1]{@#1}
\newcommand{\drp}[1]{{*}#1}
\newcommand{\nil}{\mathbf{0}}
\newcommand{\HML}{\textsf{HML}}
\newcommand{\Obs}{\mathcal{O}}
\newcommand{\Ctx}{\mathsf{Ctx}}
\newcommand{\enc}[1]{\llbracket #1 \rrbracket}
\newcommand{\dia}[1]{\langle #1 \rangle}

\title{\textbf{Inhabiting Choice}\\[0.35em]
\large Confluent Witnesses, the Phenomena They Cannot Carry,\\
and One Observation Made Three Times --- Then Turned on Its Interlocutor}

\author{L.~G.~Meredith\\ \small F1R3FLY.io}
\date{November 2026}

\begin{document}
\maketitle

\begin{abstract}
A system whose witnesses reduce confluently can \emph{state} non-confluent
phenomena but cannot \emph{inhabit} them: no term of such a system has a
branching structure, so the system's identity on its own witnesses is forced
to be linear-time, and the logic adequate for that identity is forced to be
determinate. This note makes that one observation three times. Abstractly, it
is a consequence of Church--Rosser alone: on a confluent witness language the
linear-time/branching-time spectrum collapses, whereas on the rho calculus it
does not, and the generated Hennessy--Milner logic is obliged to discriminate
proofs by the structure of their choices. For software verification, it says
that a $\lambda$-term-based prover can certify a concurrent system only
relative to a stipulated model of its non-determinism, and that the
faithfulness of that model, unlike the faithfulness of a deterministic machine
model, is not auditable by any means the prover possesses; a prover whose
witnesses are processes on a process runtime makes the certificate an
inhabitant of the system and the audit an internal bisimulation check. For
artificial intelligence, it says that branching-time structure is invisible in
any corpus of completed runs --- a trace distribution has a unique
autoregressive factorisation --- so that an intelligence trained only by
observation is linear-time by an information-theoretic limitation rather than
a design choice, while naturally evolved intelligence demonstrably exploits
races as races. The three registers share one spectrum, which we align with
Pearl's ladder, and one remedy, which is witnesses that can be races and an
identity on them at least as fine as step bisimulation. The note was drafted
in dialogue with a transformer-derived language model, and the observation is
then turned on that circumstance: the model can discuss linear versus
branching time and cannot be a witness of either, the agent loop around it
supplies intervention from outside in exactly the manner of a tactic engine
around a $\lambda$-prover, and the dialogue itself is a bisimulation game
played on the model by its interlocutor. The closing question is what would
move a culture to bear witness to its own inhabitation of that saddle point.
Nature answers it: Cohen's desert annuals inherit a race uncollapsed, the
immune system and the slime mould run races as their search procedure, and
recursive self-improvement, examined in this light, needs an improvement
operator that acts on witnesses that are races, preserves branching through
evaluation, and funds the branches that lose --- none of which an operator on
confluent outcomes can do.
\end{abstract}

% =====================================================================
\section{The observation}
% =====================================================================

The prior notes of this series~\cite{ProofSearch,PaPContinued} argued that
proof objects in Lean, Coq and Agda --- terms of a typed $\lambda$-calculus ---
force proof \emph{search} into an orchestration layer outside the proof
dynamics, because confluence leaves no room for choice in the reduction
relation and closedness leaves no open end for an incomplete proof to reach
out through. The present note isolates what we take to be the underlying fact
and states it without reference to search.

\begin{quote}
\itshape A system whose witnesses reduce confluently can talk about
non-confluent phenomena, but no witness of the system can be one. The
phenomena are present in the building and absent from the proofs.
\end{quote}

The consequence that matters is for \emph{identity}. Whatever notion of
``same proof'' such a system adopts, it is computed from a confluent
dynamics, and on a confluent dynamics the distinction between branching-time
and linear-time equivalences does not exist. The identity is therefore
linear-time whether or not anyone chose it to be, and the logic adequate for
it cannot distinguish two witnesses by \emph{when} a choice was made, because
no witness ever makes one.

We make the point three times. \S\ref{sec:abstract} states it as a pair of
propositions about witness languages, with the $\lambda$-calculus and the rho
calculus~\cite{MeredithRadestock05} as the two poles. \S\ref{sec:verification}
reads it for the verification of concurrent software, where it says precisely
what a $\lambda$-term-based prover's certificate of a concurrent system is a
certificate \emph{of}. \S\ref{sec:ai} reads it for artificial intelligence,
where it becomes a statement about what can be learned from traces and what
requires intervention. \S\ref{sec:interlocutor} turns the observation on the
circumstance of the note's own drafting, in dialogue with a transformer-derived
language model, and finds the model, the loop around it, and the dialogue
itself each sitting at a predicted rung. \S\ref{sec:saddle} asks what would
move a culture off the saddle point so located, and takes the answer from
organisms that inherit races uncollapsed, ending with recursive
self-improvement as the case where the answer is forced. \S\ref{sec:one} lays
the readings side by side, and \S\ref{sec:obligations} records what the note
leaves owed.

Throughout, the identity relation is the weak context-labelled bisimilarity of
the companion papers~\cite{BIHOL,NoWNS,Extension}: labels are minimal enabling
contexts in the open form of the Leifer--Milner
construction~\cite{LeiferMilner00}, matched up to the relation being defined,
with internal steps silent. We write $\bisW{W}$ for this relation on a witness
language $W$, and $\treq{W}$ for the corresponding weak trace equivalence.

% =====================================================================
\section{The point, abstractly}\label{sec:abstract}
% =====================================================================

\subsection{Collapse of the spectrum under confluence}

Let $W$ be an interactive graph-structured lambda theory in the sense
of~\cite{Omnibus,Extension}: a signature, an equational theory, and rewrite
rules whose left-hand sides are cuts. Its reduction relation $\to_W$ is the
closure of the rules under contexts and equations; its context-labelled
transition system has transitions $P \xrightarrow{c} P'$ for minimal enabling
contexts $c$, with the empty context labelling internal steps.

\begin{definition}[Determinate]
$W$ is \emph{determinate} when $\to_W$ is confluent and internal steps are
inert for labelled transitions: if $P \to_W P'$ and $P \xrightarrow{c} Q$
with $c$ non-empty, then $P' \xRightarrow{c} Q'$ for some $Q' \bisW{W} Q$.
\end{definition}

Confluence alone does not give inertness, but for the $\lambda$-calculus the
two come together: Church--Rosser supplies confluence, and standardisation
supplies inertness, since any labelled transition --- an application to a
fresh argument, in the open label discipline --- commutes with internal
$\beta$-steps up to further internal steps.

\begin{proposition}[Collapse]\label{prop:collapse}
If $W$ is determinate then, on terms without divergence,
$\bisW{W} \;=\; \treq{W}$. The linear-time/branching-time spectrum between
them is a single point.
\end{proposition}

\begin{proof}[Proof sketch]
This is Milner's determinacy theorem~\cite{MilnerCCS} in the labelled setting
of Groote and Sellink's $\tau$-confluence~\cite{GrooteSellink}. Trace
equivalence is contained in bisimilarity on determinate systems because the
set of states reachable by a given weak trace is, up to $\bisW{W}$, a
singleton: any two such states are internal-reducts of a common ancestor and
reconverge by confluence, and their labelled capabilities agree by inertness.
So trace-equivalent terms have bisimilar successors after every label, and
trace equivalence is a bisimulation. The reverse inclusion always holds.
Divergence is excluded because a divergent term has traces without having the
matching derivatives; on the strongly normalising fragment the hypothesis is
vacuous.
\end{proof}

\begin{corollary}[$\lambda$-witnesses]\label{cor:lambda}
Let $\Lam$ be the $\lambda$-calculus as an interactive GSLT with cut
application and rule $\beta$. On closed, typed, strongly normalising terms,
$\bisW{\Lam}$ coincides with $\beta\eta$-equality.
\end{corollary}

\begin{proof}[Proof sketch]
By Proposition~\ref{prop:collapse}, $\bisW{\Lam}$ is trace equivalence, and
with labels ``apply to a fresh argument'' traces are applicative behaviour.
Two $\beta\eta$-distinct normal forms of the same simple type are separated by
an applicative context given two base constants (the typed B\"ohm
theorem~\cite{Statman83}), and $\beta\eta$-equal terms are applicatively
indistinguishable. On this fragment $\eta$ is behaviourally invisible, since
every closed term of arrow type converges to an abstraction.
\end{proof}

So the proof identity of a $\lambda$-witnessed system is \emph{forced} to be
$\beta\eta$ on the proof-assistant fragment, and forced to be linear-time
everywhere: it was never a choice.

\subsection{Non-collapse on the rho calculus}

The rho calculus has no choice primitive; non-determinism arises from races
on names. Write $\mathsf{race}(a; P, Q)$ for
$\inp{x}{a}{P} \mid \inp{x}{a}{Q} \mid \out{a}{\nil}$, which runs exactly one
of $P$, $Q$.

\begin{proposition}[Non-collapse]\label{prop:noncollapse}
There are closed rho processes $P, Q$ with $P \treq{\rhoc} Q$ and
$P \not\bisW{\rhoc} Q$, separated by a generated Hennessy--Milner formula.
\end{proposition}

\begin{proof}
Take the process-calculus textbook pair, rendered with races:
\[
P \;=\; \inp{x}{a}{\mathsf{race}(k;\ \out{b}{\nil},\ \out{c}{\nil})},
\qquad
Q \;=\; \mathsf{race}(k;\ \inp{x}{a}{\out{b}{\nil}},\ \inp{x}{a}{\out{c}{\nil}}),
\]
with $k$ private. Both have weak traces $\{\epsilon, a, ab, ac\}$ under the
obvious observers. After $P$'s input on $a$, both $b$ and $c$ remain
reachable; after $Q$'s, exactly one does. The formula
$\dia{a}(\dia{b}\top \wedge \dia{c}\top)$, with $\dia{\ell}$ the modality for
the minimal context offering $\ell$, holds of $P$ and fails of $Q$. By the
Hennessy--Milner characterisation~\cite{HennessyMilner85} they are not
bisimilar.
\end{proof}

The pair is familiar; its interest here is what it says about witnesses. Read
$P$ and $Q$ as proof-processes of a proposition whose discharge involves
choosing between lemma $b$ and lemma $c$. $P$ is the proof that keeps both
options open past the first obligation; $Q$ is the proof that commits before
it. They are the \emph{same proof} under any linear-time identity and
\emph{different proofs} under bisimulation, and the difference is exactly when
the choice was made. No $\lambda$-term can be either of them, and no
$\lambda$-witnessed system can have a proof identity that depends on the
timing of a choice, because Proposition~\ref{prop:collapse} says the timing is
never there to depend on.

\subsection{Talked about and inhabited}

The two propositions together are the observation of \S1 in formal dress.

\begin{observation}[Statement without inhabitation]\label{obs:inhabit}
A system $W$ \emph{states} a non-confluent phenomenon when some relation
definable in $W$ exhibits it --- a transition system encoded as data, with
$P$ and $Q$ of Proposition~\ref{prop:noncollapse} as two of its states. $W$
\emph{inhabits} it when some term of $W$ exhibits it under $W$'s own
reduction. A confluent $W$ can do the first and never the second. Its identity
on witnesses, and the logic adequate for that identity, are therefore blind
to the phenomenon even while its propositions quantify over it.
\end{observation}

The generated logic makes the blindness visible. On a determinate $W$, the
generated $\HML$~\cite{Extension} is itself determinate: $\dia{c}\varphi
\wedge \dia{c}\psi$ entails $\dia{c}(\varphi \wedge \psi)$ up to internal
steps, so the modalities carry no information beyond the trace they name and
the possibility of divergence. The adequate formulae for $\lambda$-witnesses
are a degenerate fragment of the adequate formulae for rho-witnesses, in which
the connective that distinguishes $P$ from $Q$ above has nothing to
distinguish.

\begin{remark}[Which comparison is being made]
It is tempting to compare the two systems by encoding $\Lam$ into $\rhoc$ and
asking whether rho observers separate $\lambda$-terms that $\lambda$ observers
identify. The companion paper~\cite{NoWNS} supplies the discipline for that
comparison: an encoding is a pair $(F,\Phi)$ on terms and contexts, bisimilarity
is indexed by an observer class, and full abstraction holds relative to
$\operatorname{im}\Phi$ while the restriction to $\operatorname{im}\Phi$ is
forced. Carried out for $\Lam$, the comparison yields three nested relations on
$\lambda$-terms: relative to encoded $\lambda$-contexts, no separation
(Conjecture~\ref{conj:fullabs} below); relative to encoded $\pi$-contexts, the
gap between L\'evy--Longo tree equality and lazy observational
equivalence~\cite{Sangiorgi94,Lassen99}, confined to terms with unsolvable
subterms; relative to all rho contexts, a further separation that exposes the
encoding's bookkeeping and is uninformative about proofs. That comparison is
worth making, and we record its obligations in \S\ref{sec:obligations}, but it
is \emph{not} the point of this note. Every one of those relations is a
relation on $\lambda$-terms, and every $\lambda$-term is confluent. The
separation that matters is between the confluent image of $F$ and the
non-confluent witnesses outside it, where there is nothing on the $\lambda$
side to compare against. Observation~\ref{obs:inhabit} is about that.
\end{remark}

\subsection{Beyond branching: steps and causes}

The spectrum has a third level above branching time. Ordinary context-labelled
bisimulation satisfies the expansion law, $a \mid b \bis a.b + b.a$ with $+$
rendered by a race, so a rho-witnessed system with that identity inhabits
branching but not causal independence --- Observation~\ref{obs:inhabit}
transposed one level up. The remedy is again a matter of presentation rather
than of redefining bisimulation. Adding the synchronous rule
\[
\frac{P_1 \to Q_1 \qquad P_2 \to Q_2}{P_1 \mid P_2 \;\to\; Q_1 \mid Q_2}
\tag{Par-2}
\]
to $R$ as a single step --- the one-cell shadow of horizontal composition in
rewriting logic~\cite{Meseguer92} --- makes the derived transition system a
step transition system, the induced identity step bisimulation, and the
expansion law false. Step bisimulation records simultaneity and not
dependence: $a \mid b$ and $(a \mid b) + a.b$ are step bisimilar and
pomset-distinct~\cite{GlabbeekGoltz}. For dependence the rho calculus has an
unusual resource: a reduction creates at most one name, and creates it in a
substitution-transparent position~\cite{NoWNS}, so a redex on a name created
by an earlier reduction is causally after it and the quote structure of the
name is a provenance record. A history-preserving bisimulation for rho can read
its dependence relation off name creation on the fresh-name fragment; the rest
is the subject of the history-functor note.

\subsection{The price}

$\beta\eta$ on strongly normalising typed terms is decidable; bisimilarity on
a Turing-complete reflective calculus is not. The collapse of
Proposition~\ref{prop:collapse} is what buys Lean its decidable proof
identity and its cheap checking, and for propositions that are static ---
inhabitation, arithmetic --- it loses nothing, because there is nothing
behavioural to lose. The observation bites exactly where the propositions are
about behaviour: there a confluent witness cannot carry the branching
structure that \emph{is} the content of the proof, and the identity it imposes
equates proofs that differ in the one respect the proposition is about.

% =====================================================================
\section{The point, for software verification}\label{sec:verification}
% =====================================================================

\subsection{What a $\lambda$-term prover does certify}

It is important not to overstate the case. Lean, Coq and Isabelle verify
non-confluent systems routinely. Iris~\cite{Iris} proves linearisability and
race-freedom of fine-grained concurrent data structures; CertiKOS~\cite{CertiKOS}
certifies a concurrent operating-system kernel, scheduler included, down to
assembly, through a stack of certified abstraction layers and a
concurrency-aware compiler whose correctness proof is itself a Coq term. A
ticket-reservation system's ``no double booking under any interleaving'' is a
universally quantified statement over a transition relation, and its proof is a
$\lambda$-term doing induction on reachable states with an invariant. Nothing
in \S\ref{sec:abstract} prevents this: the non-confluence is \emph{stated},
as an inductive relation, and the proof is about the relation.

So the claim is not that such systems cannot be verified with $\lambda$-term
witnesses. It is a claim about what the certificate is a certificate of, and
about which part of the trusted base can be examined.

\subsection{What is outside: the model of non-determinism}

Every such development begins by stipulating a model of the non-determinism:
the interleaving semantics, the scheduler, the memory model, entered as a
definition of which executions exist. The proof quantifies over that
definition. The faithfulness of the definition to the machine is not a theorem
anywhere; it is the trusted base, as the instruction-set model is the trusted
base of a sequential verification. The asymmetry between the two cases is the
point.

\begin{observation}[Unauditability]\label{obs:audit}
For a sequential system the trusted base is a deterministic step function, and
its faithfulness is auditable by execution: run the machine, run the model,
compare traces, and every discrepancy is a counterexample found by running.
For a concurrent system the trusted base is a relation --- a set of permitted
interleavings --- and no finite set of observed runs distinguishes it from
many other relations. Faithfulness is a branching-time property of the model,
and checking it requires a branching-time observer. By
Observation~\ref{obs:inhabit}, a $\lambda$-term-based prover contains no
inhabitant with branching structure, hence nothing that can be run against the
machine to perform the check.
\end{observation}

The trust is therefore not merely unproven; it is unexaminable by any means
the prover possesses. The sequential verifier's trusted base is auditable by
the one thing a $\lambda$-prover does have, execution. The concurrent
verifier's is not.

\subsection{Cost, and where the certificate lives}

Two further consequences follow, both visible in the heroic projects.

First, cost. In the $\lambda$-witnessed world the certificate is a term checked
once and discarded, and the running system is a different object produced by
extraction or compilation. The link between them is a theorem --- a verified
compiler whose correctness statement is a simulation under all interleavings
--- and at every refinement layer the branching structure the witness cannot
carry must be re-encoded by hand as a relation and the refinement re-proved as
a $\lambda$-term \emph{about} two relations. CertiKOS reaches the bottom of its
stack at a cost proportional to exactly the structure its proof objects cannot
inhabit, and the next concurrent system starts over.

Second, deployment. A $\lambda$-proof cannot be run against the system it
certifies. A rho-witnessed proof with pending obligations \emph{is} a process,
--- under the obligation-as-communication reading of~\cite{ProofSearch} it is a
runtime monitor --- and on a runtime that is the calculus, context-labelled
bisimulation quantifies over the environments the system will actually meet,
because the contexts \emph{are} rho processes. The model/implementation gap
does not vanish; it moves, from ``this model is faithful to the machine''
(per system, unauditable) to ``this implementation is faithful to the
calculus'' (once, for the platform). That is what verified compilation did for
sequential code. Process witnesses do it for concurrent code by construction
rather than by a per-system refinement stack.

\subsection{True concurrency makes the gap wider}

The interleaving models of \S3.1 are a defensible projection on a single core
or under sequential consistency. On multicore hardware they are false, and the
field models weak memory axiomatically, as partial orders on events per
candidate execution~\cite{Herd}, because interleaving cannot express the
reorderings the hardware permits. The trusted base is then a partial-order
model, and Observation~\ref{obs:audit} sharpens: its faithfulness is a property
of an independence relation, two levels above anything a confluent witness can
inhabit. On the rho side the corresponding move is Par-2 of \S2.4; the
F1R3FLY runtime is itself a partial-order system, a block DAG merging
independent deploys by conflict detection, so a verification on that runtime
whose identity is interleaving bisimulation under-specifies the system it runs
on. For the ticket system the difference is concrete: two reservations on
disjoint seats should commute, an interleaving proof reaches that fact only
through a linearisability argument that re-derives commutation from outside,
and a step-concurrent witness has the commutation as a single Par-2 step.

% =====================================================================
\section{The point, for artificial intelligence}\label{sec:ai}
% =====================================================================

\subsection{Traces cannot reveal branching}

The value proposition of the current programme is that the salient aspects of
intelligence are software, and in its dominant form, that they are what one
obtains by predicting traces at scale. Proposition~\ref{prop:noncollapse}
has a direct consequence for that form.

\begin{proposition}[Observational invisibility of branching]\label{prop:invisible}
Let $\mu$ be a probability distribution over finite traces. Then $\mu$ has a
unique autoregressive factorisation $\mu(x_1\cdots x_n) = \prod_i
\mu(x_i \mid x_{<i})$, and any two labelled transition systems inducing $\mu$
--- in particular the probabilistic renderings
$p\cdot a.b + (1-p)\cdot a.c$ and $a.(p\cdot b + (1-p)\cdot c)$ of
Proposition~\ref{prop:noncollapse} --- induce the same factorisation. A
learner whose only access to the system is samples from $\mu$ cannot
distinguish them.
\end{proposition}

\begin{proof}
Conditional probabilities are determined by the joint distribution; the two systems are
trace-distribution equivalent by construction and not probabilistically
bisimilar~\cite{LarsenSkou91}.
\end{proof}

So a next-token predictor trained on a corpus of completed runs, however
large, represents the trace distribution and nothing more. The question
``was the die cast before or after $a$?'' has no answer in the data. The
limitation is not architectural; it is information-theoretic.
\emph{Branching-time structure is invisible to any observer that does not
act.}

\subsection{Experimenters and spectators}

To learn branching structure one has to intervene: place the system in a
state, supply a context, and observe what it can still do. That is the
context-labelled bisimulation game of~\cite{BIHOL,Extension} read as an
epistemology --- the observer chooses a label, the system answers --- and the
mortal-scientist programme~\cite{Scientist} is its agent. Bisimulation is the
equivalence an \emph{experimenter} can learn; trace equivalence is what a
\emph{spectator} gets.

The alignment with Pearl's ladder~\cite{Pearl09} is exact. Rung one,
association, is linear time: what a trace distribution supports. Rung two,
intervention, is branching time: what an experimenter supplying contexts can
learn, and what bisimulation characterises. Rung three, counterfactuals,
requires the causal structure, which is the true-concurrency level of
\S2.4 --- history-preserving bisimulation, with dependence as an object. The
ladder and the linear/branching/causal spectrum are one hierarchy seen from
statistics and from concurrency theory.

\subsection{Where the stack sits}

The indictment is specific and should be aimed carefully. Next-token
prediction on observational text is rung one by construction. Reinforcement
learning is not: a Markov decision process is a branching-time object, a
state with actions and distributions over successors, and bisimulation for
MDPs is an established tool for state abstraction~\cite{GivanDeanGreig03}. An
RL agent learns by intervening. Post-training a language model with outcome
rewards is the community purchasing rung-two data for a rung-one model, which
is evidence that the limitation is felt.

But reinforcement learning stops one rung short of nature. An MDP, and even a
Markov game with simultaneous moves, has a global clock: the world steps in
lockstep. The world does not. The races of the natural world --- which sperm
fertilises the egg, which antelope the pride culls, which flower the bee
reaches --- are truly concurrent, and an agent that models them as
interleaved under a global tick attributes dependence where there is only
coincidence. That is the step/pomset gap of \S2.4, now as a claim about what
evolved intelligence must track.

\subsection{Nature exploits races as races}

The examples must be of the right kind. Sensitivity to \emph{who won} a race
is observation of an outcome, rung one. The claim requires organisms sensitive
to the race \emph{as a race} --- that it could go otherwise, and that one can
act on the branching. The ethological and evolutionary literature supplies
them.

\begin{itemize}[leftmargin=2.2em,itemsep=0.25em]
\item \textbf{Protean behaviour.} Prey move unpredictably in order to defeat
a predator's predictive model~\cite{HumphriesDriver70}. The strategy only
makes sense against an opponent reasoning in branching time, and only works if
the prey itself can generate genuine branching.
\item \textbf{Mixed strategies.} Hawk--dove and every mixed evolutionarily
stable strategy~\cite{MaynardSmith82} is a population exploiting the fact that
the opponent must model a choice, not a trace.
\item \textbf{Bet-hedging.} Organisms produce offspring carrying non-confluent
strategies, so that whichever environment arrives some
win~\cite{SegerBrockmann87}. This is the biological hedge: pursue both
branches, keep what succeeds.
\item \textbf{Sexual selection.} The race is manufactured; the ornament exists
to structure a competition whose outcome is informative only because it was a
race.
\end{itemize}

The lion that drives the herd to \emph{create} a race it can win is doing
rung-two work; the lion that eats the slowest is doing rung one. The first is
intelligence; the second is a filter.

\subsection{Witnesses, again}

The last step returns to \S\ref{sec:abstract}. An agent's knowledge is its
hypotheses and plans --- in the mortal-scientist setting, its proofs.
Bet-hedging is a \emph{plan that is itself a race}: pursue two branches in
parallel and commit to whichever lands. That plan is a non-confluent witness.
A $\lambda$-plan cannot be it; it can describe it from outside and hand it to
a scheduler, which is the orchestration-outside-the-dynamics of~\cite{ProofSearch}
reappearing as a limit on cognition. So the requirement on a software
rendering of intelligence is not only that it \emph{discriminate} branching
phenomena but that its internal representations \emph{inhabit} them.
Observation~\ref{obs:inhabit}, drawn through the agent's own cognition: an
intelligence whose hypotheses are confluent can observe hedging but cannot
hedge in its hypotheses.

\subsection{The smallest behaviour}\label{sec:smallest}

It is worth saying how small the stymieing instance is, because the
smallness is the argument's strength. Milner's vending
machines~\cite{MilnerCCS},
\[
P \;=\; \mathit{coin}.(\mathit{tea} + \mathit{coffee})
\qquad\qquad
Q \;=\; \mathit{coin}.\mathit{tea} + \mathit{coin}.\mathit{coffee},
\]
are the pair of Proposition~\ref{prop:noncollapse} with three actions, and
there is no trace-equivalent, non-bisimilar pair with fewer. Every completed
session of either machine is $\mathit{coin}\;\mathit{tea}$ or
$\mathit{coin}\;\mathit{coffee}$, at identical frequencies if $Q$ flips the
coin its customer would. The observable a trace-trained learner cannot
recover is a \emph{refusal}: coin in, tea requested, nothing out. $Q$ can
refuse tea after a coin; $P$ cannot; and a refusal never appears in a log of
completed runs, because a log records what happened and not what could not
have happened. Text records what was said, never what could not have been
said.

The test is correspondingly small. Generate sessions from one of the two
machines chosen by a hidden coin; record completed actions only; train, or
prompt in context; ask whether, after a coin, the machine can refuse tea.
Chance performance is forced rather than likely, since the two training
distributions are one distribution. The remedy is equally small: let the
learner \emph{try}. One attempted tea after one coin, met with a refusal,
settles the question, and no quantity of further logs can substitute for
it. That single attempt is the step from spectator to experimenter, from
Pearl's first rung to his second, and from trace equivalence to
bisimulation. Two fairness remarks: a learner handed the \emph{term} $Q$ in
its prompt answers correctly, since the stymie is for learning from
observation and not for reasoning about a description; and logs that record
failed attempts make the refusal a trace symbol and the distinction
learnable, but such logs are intervention data in disguise, which is the
point rather than a loophole. That the same three-action pair is the smallest
behaviour invisible to a $\lambda$-prover's proof identity
(Proposition~\ref{prop:noncollapse}) is the note's thesis in one example.

% =====================================================================
\section{The point, for the interlocutor}\label{sec:interlocutor}
% =====================================================================

This note was drafted in conversation with a transformer-derived language
model, which read the companion papers, proposed the propositions of
\S\ref{sec:abstract}, conceded the CertiKOS counterexample of
\S\ref{sec:verification}, and compiled the document. A reader is entitled
to ask how a trace-trained system can have had that conversation if
\S\ref{sec:ai} is right. The answer is the observation of \S1 applied to its
own author's assistant, and it comes in three refinements.

\subsection{Statement, not inhabitation}

The thesis never said a transformer cannot \emph{talk about} branching time.
Bisimulation, refusals, the vending machines, Pearl's ladder are in the corpus
as described objects, and reasoning about descriptions is what a trace-trained
system does. The model's contribution to this note is Lean formalising the
$\pi$-calculus: a relation stated, theorems proved about it, nothing
inhabited. Apply the note's own criterion. Can the model refuse? A pure
autoregressive sampler never does; it always emits a next token, and its
choices form a fixed tree of conditionals determined by the trace
distribution, with no state an environment can probe by offering tea and
being declined. In the precise sense of \S\ref{sec:smallest}, the model is
nearer to $Q$ as described by $P$'s logs than to either machine. The
conversation is an instance of the thesis, not a counterexample to it.

\subsection{The loop}

The session was not pure sampling. Asked to read a repository, the model
cloned it; asked for a PDF, it ran the typesetter and read the log, and the
first compilation failed on an undefined control sequence --- a refusal from
the environment that no quantity of training text would have supplied, and
that was corrected on the next attempt. The harness around the model is an
agent loop, and an agent loop is Pearl's second rung. So what the author was
in conversation with is a rung-one model inside a rung-two loop: the same
architecture as a $\lambda$-prover inside a tactic engine, with the
orchestration outside the dynamics once more. The critique of
\S\ref{sec:ai} applies to the model; the loop is the partial remedy; and the
thesis says the remedy is partial because the branching lives in the harness
rather than in the witnesses, which is the subject of \S\ref{sec:saddle}.

\subsection{The dialogue}

The third refinement concerns the reader rather than the model. The
conversation that produced this note was a bisimulation game with the author
as experimenter. She offered contexts --- ``does your example not confirm my
argument?'', ``does this hold for true concurrency?'' --- and learned where
the model concedes, where it pushes back, and where it corrects itself. That
is rung-two knowledge about the model, obtained by intervention, that no
transcript of its previous conversations would have given. The observer
extension of~\cite{Extension} says that structural facts become behavioural
once the observer holds instruments for taking a term apart; a pointed
question is such an instrument, and the author held several. The note's
point, demonstrated on the note's assistant, by the note's author.

% =====================================================================
\section{The saddle point}\label{sec:saddle}
% =====================================================================

\S\ref{sec:interlocutor} located the present culture's working arrangement:
a confluent predictor whose hedges --- best-of-$n$, tree-of-thought,
sample-and-verify --- live in a harness. The arrangement works, and that is
the difficulty. A harness supplies the second rung, which is enough to be
useful; so the question is Socratic rather than technical. \emph{What moves a
culture to bear witness to its own inhabitation of a saddle point?} The
answer offered here is that nothing moves it while it is buying answers,
and that one thing moves it the moment it wants inheritance. Nature supplies
both the diagnosis and the smallest instance.

\subsection{A hedge that is inherited}

Cohen's desert annuals~\cite{Cohen66}: in a year of uncertain rain, a
fraction of a plant's seeds germinate and the rest lie dormant. The fraction
is heritable, and selection tunes it to the variance of the environment. The
organism is not predicting which year it is in. It is running a race between
two branches --- germinate now, wait --- and \emph{passing the race on} to the
next generation as a race, encoded in the genome, uncollapsed.

That is the whole mechanism in one organism. The witness carries the
branching; the branching is inherited; selection acts on the branching
structure itself. Fisher's theorem needs heritable variation in the
\emph{strategy}, and a hedge is evolvable only if the hedge is in the heritable
object. Protean escape, mixed strategies, sperm competition are downstream of
the same fact: the capacity to generate a race is in the witness, so it can
be selected on.

\subsection{Why the saddle point is a saddle point}

Now ask where the loop architecture keeps its race. Best-of-$n$ is a hedge,
and it works; but the hedge lives in the harness. The harness runs $n$
samples, picks a winner, and emits a string. The race is collapsed at the
boundary, paid for again on every query, and handed to no one as a race. It
is not heritable. Nothing downstream can select on it, compose it, or run it
in a different environment, because downstream receives an outcome and not a
structure.

The loop architecture is the desert annual that germinates every seed every
year and reseeds from whichever survived. It survives, and it cannot
\emph{accumulate}. It answers; it does not do science in the sense
of~\cite{Scientist}: hypotheses as objects that carry their own branching, are
tested by environments the hypothesiser did not choose, and are inherited by
the next hypothesiser with the branching intact. The saddle point is stable
exactly as long as per-query answers are what is being paid for. The moment
the thing paid for is an accumulating body of reusable, selectable, composable
search structure --- a scientific community rather than an oracle --- the
per-query collapse is the dominant cost, and the race has to move from the
harness into the witness.

\subsection{Organisms that are the mechanism}

Some natural systems do not merely exploit a race; they \emph{are} the search
mechanism of~\cite{ProofSearch}, and they show the design built and tested at
scale.

\begin{itemize}[leftmargin=2.2em,itemsep=0.25em]
\item \textbf{Clonal selection}~\cite{Burnet59}. Antibodies are hypotheses,
antigen binding is the obligation, lymphocytes race to bind, winners
replicate, losers starve. Massively parallel, non-confluent, coordinator-free
proof search with mortality as pruning and a proven lemma --- a memory cell
--- as a reusable metabolite.
\item \textbf{The slime mould in a maze}~\cite{Nakagaki00}. Explore every
branch in parallel, starve the unproductive tubes, keep the shortest: the
chemical-abstract-machine soup with starvation-deadlock as garbage
collection.
\item \textbf{Exuberance and pruning} in neural development, \textbf{quorum
sensing} in the swarm, \textbf{kinetic proofreading} at the molecular scale
--- each a race run \emph{as} the computation, not simulated by a scheduler
beside it.
\end{itemize}

None keeps its race in a harness, and none could have evolved if it had.

\subsection{Recursive self-improvement needs heritable races}

The seed is the challenge of recursive self-improvement in a nutshell, and
specifically the challenge rather than only the picture. Self-improvement
needs an improvement operator that acts on the thing that generates
behaviour, and the seed shows what the operator has to preserve. The
germination fraction is improvable only because the plant does not know which
year it is in: the information that would say whether the hedge was right
arrives after it has committed, often after it is dead. So the hedge must be
kept \emph{uncollapsed through the evaluation}, and the signal must land on
the next generation's copy of the hedge, not on this generation's outcome.
Improvement acts on the race. It cannot act on the result of the race,
because by the time there is a result there is nothing left to improve.

That is where the loop architecture fails as a self-improver, and the failure
has been measured. Its improvement operator is a gradient on outcomes: sample
$n$, reward the winner, move the policy toward it. Each step collapses the
race it was run on; the structure the \emph{next} step would need --- the
hedge, the diversity, the dormant seeds --- is what the current step consumes.
Loss of output diversity under outcome-based fine-tuning~\cite{Kirk24} is the
symptom; the cause is that the operator acts on results, and results are
confluent. An operator on confluent objects is self-limiting by
construction: it climbs, and eats the branches it would need to climb
further.

Nature's operator lacks the defect because it acts on the witness and pays
for the losers. Dormant seeds are branches that did not win this year, funded
anyway, and the fraction funded is itself the heritable, improvable
parameter. In the terms of this series it is the phlogiston economy with a
non-zero budget for losing branches, the budget part of the genome:
selection tunes how much of the search's energy goes to hedges, and that
tuning is the improvement.

\begin{observation}[The three requirements]\label{obs:rsi}
An improvement operator capable of recursion must (i) be defined on witnesses
that are races, (ii) preserve branching through evaluation, and (iii) fund
branches that lose at a rate that is itself subject to the operator. An
operator on confluent outcomes satisfies none of the three.
\end{observation}

A loop around a confluent predictor can meet the three requirements only by
rebuilding the mechanism outside the model, which is the saddle point. A
reflective calculus with an internal cost economy meets them natively:
tactics are processes, processes are quotable data, and the economy funds
them at a rate that is itself a process. The plant does this with one
parameter. The question for the programme is whether it can be done with the
whole genome of rewrite rules, which is the variation operator
\cite{ProofSearch} already names.

\subsection{The honest version}

A determined harness engineer can make the loop emit a portfolio rather than
a winner: output $k$ candidates and let the consumer choose. But then the
consumer must be able to \emph{run} the portfolio under its own environment,
which makes the consumer a process runtime and the artefact passed between
them a process. At that point the rho mechanism has been rebuilt at the
boundary, and the question is only whether to build it once as the substrate
or repeatedly at every interface. So there is no phenomenon that forces a
culture off the saddle point while it is buying answers. There is one that
forces it the moment it wants inheritance: a hedge that has to survive its
author. Nature found that phenomenon first, and the fraction of seeds that
stay in the ground is its smallest instance.

% =====================================================================
\section{One point, three registers}\label{sec:one}
% =====================================================================

\begin{table}[h]
\centering
\small
\begin{tabular}{@{}p{0.2\textwidth}p{0.25\textwidth}p{0.25\textwidth}p{0.22\textwidth}@{}}
\toprule
 & \textbf{abstract} & \textbf{verification} & \textbf{intelligence} \\
\midrule
the confluent object & a $\lambda$-witness & a proof about a model & a trace predictor \\
what it can state & a transition relation & the system's interleavings & the world's runs \\
what it cannot inhabit & a race & the system's non-determinism & a hedge \\
forced identity & $\beta\eta$ / trace eq. & model-relative certificate & trace distribution \\
what is unexaminable & timing of choice & faithfulness of the model & when the die was cast \\
the remedy & rho witnesses, $\bisW{\rhoc}$ & process certificates on a process runtime & an experimenter with race-shaped hypotheses \\
smallest instance & $a.(b+c)$ vs $a.b+a.c$ & a refusable request after a commit & the vending machine, logs only \\
the saddle point & orchestrator beside the dynamics & tactic engine around the prover & harness around the predictor \\
what leaves it & tactics as witnesses & certificates that run & hedges that are inherited \\
one level up & Par-2, provenance & weak memory, block DAGs & clockless nature, rung three \\
\bottomrule
\end{tabular}
\caption{The observation in its three readings.}
\label{tab:three}
\end{table}

The rows of Table~\ref{tab:three} are not analogies. The same proposition ---
collapse of the spectrum under confluence --- underlies all three columns, and
the same remedy answers all three: witnesses that can be races, and an identity
on them at least as fine as step bisimulation, with the generated logic
following the presentation. The verification column and the intelligence
column are the abstract column applied to a certificate and to a hypothesis
respectively; a certificate is a hypothesis one has stopped testing. The
saddle-point row is the same arrangement three times --- confluent object,
branching supplied from outside --- and the row beneath it is the same
departure three times: put the race in the object that is kept.

% =====================================================================
\section{Obligations}\label{sec:obligations}
% =====================================================================

The note is an observation with consequences, and it leaves theorems owed.

\begin{obligation}[Determinacy of $\Lam$]
Establish inertness for the open minimal-label transition system of $\Lam$
under the BIHOL construction, so that Proposition~\ref{prop:collapse} applies
without the standardisation hand-wave, and identify $\bisW{\Lam}$ exactly on
untyped terms: lazy observational equivalence if the labels are closed
second-order contexts~\cite{DHL09}, normal-form bisimilarity --- L\'evy--Longo
tree equality~\cite{Lassen99} --- if they are open with fresh variables. The
two differ on terms with unsolvable subterms and nowhere else.
\end{obligation}

\begin{conjecture}[Full abstraction relative to encoded contexts]\label{conj:fullabs}
For the composite encoding $\enc{-} : \Lam \to \rhoc$ through Milner's
encoding~\cite{Milner92} and the local-allocator encoding of~\cite{NoWNS},
$M \bisW{\Lam} N$ iff $\enc{M} \bisW{\rhoc,\operatorname{im}\Phi} \enc{N}$.
The $\pi$-to-rho half is the full-abstraction theorem of~\cite{NoWNS}; the
$\lambda$-to-$\pi$ half is operational correspondence for the standard
encoding, with the base case of soundness the place to watch.
\end{conjecture}

\begin{obligation}[Step and causal identities]
Verify that Par-2 as a rule scheme admits redex-RPOs in the sense
of~\cite{BIHOL}, so that step labels are multisets of minimal contexts and
step bisimulation is the derived identity; then align the provenance relation
read off name creation with history-preserving bisimulation on the fresh-name
fragment, and state what the history functor adds on the remainder.
\end{obligation}

\begin{obligation}[The audit as a theorem]
Observation~\ref{obs:audit} is stated informally. Its content is that
faithfulness of a non-deterministic model is not a trace property, which is
Proposition~\ref{prop:invisible} read for machines rather than corpora; a
precise statement should say which class of tests a $\lambda$-prover can run
against a system and prove that the class is contained in trace observation.
\end{obligation}

\begin{obligation}[The vending-machine benchmark]
Run the experiment of \S\ref{sec:smallest}: logs from one of $P$, $Q$ chosen
by a hidden coin, completed actions only, and the refusal question. The
prediction is chance for any learner on the logs and one query for a learner
permitted to attempt; the experiment is cheap and its failure would be
informative, since it would mean refusals leak into completed-run logs by a
channel the argument has not named.
\end{obligation}

\begin{obligation}[The three requirements as a theorem]
Observation~\ref{obs:rsi} is stated informally. Its content should be a
statement about improvement operators on a witness language $W$: if $W$ is
determinate then every operator definable on $W$-outcomes is invariant under
$\treq{W}$ and so cannot increase the branching of what it improves, whereas
on $\rhoc$ with the cost endofunctor the funding rate of losing branches is
itself a process and hence in the operator's domain. The tactic-evolution
target of~\cite{ProofSearch} is the constructive half.
\end{obligation}

\begin{obligation}[Target 3 of~\cite{ProofSearch}, relocated]
``Proof identity is bisimulation'' is trivial on saturated proof-processes,
which have no transitions and are therefore all bisimilar to $\nil$. The
target should be restated on open proof-processes --- proofs as functions of
their obligations --- where \S2.1 says it agrees with $\beta\eta$ on the
$\lambda$-image and \S2.2 says it does not collapse outside it; exported
witnesses, compared as names, are identified up to structural congruence and
belong to a separate proposition.
\end{obligation}

\section*{Acknowledgements}

The questions that shaped this note --- whether $\lambda$-witnessed systems
are even cognizant of non-confluent phenomena, what the functional difference
is for a verifier armed only with such a prover, whether the argument survives
true concurrency, how to reconcile it with the circumstance of its own
drafting, and what moves a culture off a saddle point it has not noticed ---
are answered here in the order they were asked. Claude (Anthropic) assisted
with source review, drafting, and typesetting, and is the model of
\S\ref{sec:interlocutor}; the characterisation of its situation there was
offered by the model and is reproduced with its agreement.

\begin{thebibliography}{99}

\bibitem{ProofSearch} L.~G.~Meredith.
\newblock Proof search as concurrent reaction: the rho calculus and the
category of continued interactive GSLTs as a substrate for proof objects with
internalized, cost-accounted search.
\newblock F1R3FLY.io technical note, 2026.

\bibitem{PaPContinued} L.~G.~Meredith.
\newblock Proofs as processes, continued.
\newblock F1R3FLY.io technical note, 2026.

\bibitem{BIHOL} C.~B.~Wells, M.~Stay and L.~G.~Meredith.
\newblock Behavior in higher-order languages: admissible observers, relative
pushouts, and the encoding of extended theories.
\newblock In preparation, F1R3FLY.io, 2026.

\bibitem{NoWNS} L.~G.~Meredith.
\newblock Namespaces without a name server: freshness by communication, and a
fully abstract encoding of the $\pi$-calculus in the rho calculus.
\newblock F1R3FLY.io technical note, v3, 2026.

\bibitem{Extension} L.~G.~Meredith.
\newblock Spatial formulae, discrimination power, and the observer extension
of an interactive GSLT.
\newblock F1R3FLY.io technical note, v2, 2026.

\bibitem{Omnibus} L.~G.~Meredith.
\newblock Graph-structured lambda theories: a reference for the working
developer.
\newblock F1R3FLY.io, 2026.

\bibitem{Scientist} L.~G.~Meredith.
\newblock The mortal scientist: learning as cost-accounted computation in a
reflective higher-order calculus.
\newblock F1R3FLY.io, 2026.

\bibitem{MeredithRadestock05} L.~G.~Meredith and M.~Radestock.
\newblock A reflective higher-order calculus.
\newblock \emph{Electronic Notes in Theoretical Computer Science}
\textbf{141}(5):49--67, 2005.

\bibitem{LeiferMilner00} J.~J.~Leifer and R.~Milner.
\newblock Deriving bisimulation congruences for reactive systems.
\newblock In \emph{CONCUR 2000}, LNCS \textbf{1877}, pp.~243--258. Springer, 2000.

\bibitem{MilnerCCS} R.~Milner.
\newblock \emph{Communication and Concurrency}.
\newblock Prentice Hall, 1989.

\bibitem{GrooteSellink} J.~F.~Groote and M.~P.~A.~Sellink.
\newblock Confluence for process verification.
\newblock \emph{Theoretical Computer Science} \textbf{170}(1--2):47--81, 1996.

\bibitem{Statman83} R.~Statman.
\newblock $\lambda$-definable functionals and $\beta\eta$ conversion.
\newblock \emph{Archiv f\"ur mathematische Logik und Grundlagenforschung}
\textbf{23}:21--26, 1983.

\bibitem{HennessyMilner85} M.~Hennessy and R.~Milner.
\newblock Algebraic laws for nondeterminism and concurrency.
\newblock \emph{Journal of the ACM} \textbf{32}(1):137--161, 1985.

\bibitem{Sangiorgi94} D.~Sangiorgi.
\newblock The lazy lambda calculus in a concurrency scenario.
\newblock \emph{Information and Computation} \textbf{111}(1):120--153, 1994.

\bibitem{Lassen99} S.~B.~Lassen.
\newblock Bisimulation in untyped lambda calculus: B\"ohm trees and
bisimulation up to context.
\newblock \emph{Electronic Notes in Theoretical Computer Science}
\textbf{20}:346--374, 1999.

\bibitem{DHL09} P.~Di Gianantonio, F.~Honsell and M.~Lenisa.
\newblock RPO, second-order contexts, and $\lambda$-calculus.
\newblock \emph{Logical Methods in Computer Science} \textbf{5}(3), 2009.

\bibitem{Milner92} R.~Milner.
\newblock Functions as processes.
\newblock \emph{Mathematical Structures in Computer Science}
\textbf{2}(2):119--141, 1992.

\bibitem{Meseguer92} J.~Meseguer.
\newblock Conditional rewriting logic as a unified model of concurrency.
\newblock \emph{Theoretical Computer Science} \textbf{96}(1):73--155, 1992.

\bibitem{GlabbeekGoltz} R.~J.~van Glabbeek and U.~Goltz.
\newblock Refinement of actions and equivalence notions for concurrent systems.
\newblock \emph{Acta Informatica} \textbf{37}:229--327, 2001.

\bibitem{Iris} R.~Jung, R.~Krebbers, J.-H.~Jourdan, A.~Bizjak, L.~Birkedal and
D.~Dreyer.
\newblock Iris from the ground up: a modular foundation for higher-order
concurrent separation logic.
\newblock \emph{Journal of Functional Programming} \textbf{28}, 2018.

\bibitem{CertiKOS} R.~Gu, Z.~Shao, H.~Chen, X.~Wu, J.~Kim, V.~Sj\"oberg and
D.~Costanzo.
\newblock CertiKOS: an extensible architecture for building certified
concurrent OS kernels.
\newblock In \emph{OSDI 2016}, pp.~653--669. USENIX, 2016.

\bibitem{Herd} J.~Alglave, L.~Maranget and M.~Tautschnig.
\newblock Herding cats: modelling, simulation, testing, and data mining for
weak memory.
\newblock \emph{ACM Transactions on Programming Languages and Systems}
\textbf{36}(2), 2014.

\bibitem{LarsenSkou91} K.~G.~Larsen and A.~Skou.
\newblock Bisimulation through probabilistic testing.
\newblock \emph{Information and Computation} \textbf{94}(1):1--28, 1991.

\bibitem{Pearl09} J.~Pearl.
\newblock \emph{Causality: Models, Reasoning, and Inference}, 2nd ed.
\newblock Cambridge University Press, 2009.

\bibitem{GivanDeanGreig03} R.~Givan, T.~Dean and M.~Greig.
\newblock Equivalence notions and model minimization in Markov decision
processes.
\newblock \emph{Artificial Intelligence} \textbf{147}(1--2):163--223, 2003.

\bibitem{HumphriesDriver70} D.~A.~Humphries and P.~M.~Driver.
\newblock Protean defence by prey animals.
\newblock \emph{Oecologia} \textbf{5}:285--302, 1970.

\bibitem{MaynardSmith82} J.~Maynard Smith.
\newblock \emph{Evolution and the Theory of Games}.
\newblock Cambridge University Press, 1982.

\bibitem{SegerBrockmann87} J.~Seger and H.~J.~Brockmann.
\newblock What is bet-hedging?
\newblock \emph{Oxford Surveys in Evolutionary Biology} \textbf{4}:182--211, 1987.

\bibitem{Cohen66} D.~Cohen.
\newblock Optimizing reproduction in a randomly varying environment.
\newblock \emph{Journal of Theoretical Biology} \textbf{12}(1):119--129, 1966.

\bibitem{Burnet59} F.~M.~Burnet.
\newblock \emph{The Clonal Selection Theory of Acquired Immunity}.
\newblock Cambridge University Press, 1959.

\bibitem{Nakagaki00} T.~Nakagaki, H.~Yamada and \'A.~T\'oth.
\newblock Maze-solving by an amoeboid organism.
\newblock \emph{Nature} \textbf{407}:470, 2000.

\bibitem{Kirk24} R.~Kirk, I.~Mediratta, C.~Nalmpantis, J.~Luketina,
E.~Hambro, E.~Grefenstette and R.~Raileanu.
\newblock Understanding the effects of RLHF on LLM generalisation and
diversity.
\newblock In \emph{ICLR 2024}.

\end{thebibliography}

\end{document}
